\chapter{Boolean Analysis — Encoding Properties}
%%% generated by the blueprint pipeline from TCSlib.BooleanAnalysis.Switching.EncodingProperties
\section{Overview}

This module collects the stability, length-bound, and independence lemmas used in
Razborov's encoding argument for the switching lemma. They describe the behaviour of
the low-level routines \texttt{SwitchingLemma2.processClauseLits},
\texttt{SwitchingLemma2.razborovEncode}, and \texttt{SwitchingLemma2.razborovDecode} on
restrictions $\rho_0$, $\sigma$, decision-tree paths, and clause literals, culminating
in the technical facts needed for the encode/decode round-trip.

%%%%%%%%%%%%%%%%%%%%%%%%%%%%%%%%%%%%%%%%%%%%%%%%%%%%%%%%%%%%%%%%%%%%%%%%%%%%%%%
\section{Declarations}

\begin{lemma}[Length bound on the clause-processing output]
\lean{SwitchingLemma2.processClauseLits_bound}
\proofsource{switching-lemma}{pdf-b5e074215b9e-p002-b004}
\label{SwitchingLemma2.processClauseLits_bound}
\leanok
\proofstep
  {BoolCircuit.Lit}
  {context}
  {switching-lemma}
  {pdf-b5e074215b9e-p001-b004}
\proofstep
  {BoolCircuit.Lit.eval}
  {context}
  {switching-lemma}
  {pdf-b5e074215b9e-p001-b004}
\proofstep
  {Literal}
  {context}
  {switching-lemma}
  {pdf-b5e074215b9e-p001-b004}
\proofstep
  {SwitchingLemma2.Restriction}
  {context}
  {switching-lemma}
  {pdf-b5e074215b9e-p001-b007}
\proofstep
  {SwitchingLemma2.Literal.fixedBy}
  {context}
  {switching-lemma}
  {pdf-b5e074215b9e-p002-b004}
\proofstep
  {SwitchingLemma2.processClauseLits}
  {context}
  {switching-lemma}
  {pdf-b5e074215b9e-p002-b003,
    pdf-b5e074215b9e-p002-b002}
\uses{Literal, SwitchingLemma2.Restriction, SwitchingLemma2.processClauseLits}
\difficulty{4}
Fix a list of literal–position pairs $L$, a path $P$ (a list of variable–direction
pairs), and two restrictions $\rho_0, \sigma$ on $n$ variables. Running the
clause-processing operation on these inputs yields a remaining path $P'$ and an
auxiliary list $A$ of position–direction pairs (together with updated restrictions).
Then
\[
\abs{A} + 2\,\abs{P'} \le 2\,\abs{P},
\]
where $\abs{\cdot}$ denotes the number of entries of a list.
\end{lemma}

\begin{lemma}[Tight length bound for processing a nonempty clause against a nonempty path]
\lean{SwitchingLemma2.processClauseLits_tight}
\label{SwitchingLemma2.processClauseLits_tight}
\leanok
\uses{Literal, SwitchingLemma2.Restriction, SwitchingLemma2.processClauseLits, SwitchingLemma2.processClauseLits_bound}
\difficulty{3}
Run the clause-processing procedure on a nonempty list of free literals — each paired
with its position index in the clause — and a nonempty canonical path, starting from
restrictions $\rho_0$ and $\sigma$, and let its output consist of a remaining path and a
list of accumulated auxiliary entries (among the updated restrictions). Writing $a$ for
the number of auxiliary entries produced, $r$ for the length of the remaining path
returned, and $m$ for the length of the input path, we have
\[
a + 1 + 2r \le 2m.
\]
\end{lemma}

\begin{lemma}[Length bound for the switching-lemma encoder loop]
\lean{SwitchingLemma2.encode_go_aux_length_bound}
\proofsource{switching-lemma}{
  pdf-b5e074215b9e-p002-b004,
  pdf-b5e074215b9e-p002-b003
}
\label{SwitchingLemma2.encode_go_aux_length_bound}
\leanok
\proofstep
  {BoolCircuit.Lit}
  {context}
  {switching-lemma}
  {pdf-b5e074215b9e-p001-b004}
\proofstep
  {BoolCircuit.Lit.eval}
  {context}
  {switching-lemma}
  {pdf-b5e074215b9e-p001-b004}
\proofstep
  {Literal}
  {context}
  {switching-lemma}
  {pdf-b5e074215b9e-p001-b004}
\proofstep
  {Term}
  {context}
  {switching-lemma}
  {pdf-b5e074215b9e-p001-b004}
\proofstep
  {DNF}
  {context}
  {switching-lemma}
  {pdf-b5e074215b9e-p001-b004}
\proofstep
  {DecisionTree}
  {context}
  {switching-lemma}
  {pdf-b5e074215b9e-p001-b005}
\proofstep
  {DecisionTree.depth}
  {context}
  {switching-lemma}
  {pdf-b5e074215b9e-p001-b005}
\proofstep
  {DecisionTree.deepPath}
  {context}
  {switching-lemma}
  {pdf-b5e074215b9e-p002-b003}
\proofstep
  {SwitchingLemma2.Restriction}
  {context}
  {switching-lemma}
  {pdf-b5e074215b9e-p001-b007}
\proofstep
  {SwitchingLemma2.Restriction.freeVars}
  {context}
  {switching-lemma}
  {pdf-b5e074215b9e-p001-b007}
\proofstep
  {SwitchingLemma2.Restriction.numFree}
  {context}
  {switching-lemma}
  {pdf-b5e074215b9e-p001-b007}
\proofstep
  {SwitchingLemma2.Literal.fixedBy}
  {context}
  {switching-lemma}
  {pdf-b5e074215b9e-p002-b001}
\proofstep
  {SwitchingLemma2.canonicalDTree}
  {context}
  {switching-lemma}
  {pdf-b5e074215b9e-p002-b002}
\proofstep
  {SwitchingLemma2.processClauseLits}
  {context}
  {switching-lemma}
  {pdf-b5e074215b9e-p002-b003}
\proofstep
  {SwitchingLemma2.razborovEncode}
  {context}
  {switching-lemma}
  {pdf-b5e074215b9e-p002-b003,
    pdf-b5e074215b9e-p002-b004}
\proofstep
  {SwitchingLemma2.processClauseLits_bound}
  {context}
  {switching-lemma}
  {pdf-b5e074215b9e-p002-b004}
\proofstep
  {SwitchingLemma2.processClauseLits_tight}
  {context}
  {switching-lemma}
  {pdf-b5e074215b9e-p002-b004}
\uses{DNF, SwitchingLemma2.Restriction, SwitchingLemma2.processClauseLits_tight}
\difficulty{4}
Consider the recursive encoding routine underlying Razborov's proof of the switching
lemma, which processes a DNF formula $f$ on $n$ variables against a width parameter $w$
and a fuel bound, while carrying a path recorded as a list of variable–value pairs
$(i,b) \in \mathrm{Fin}\,n \times \mathrm{Bool}$, two restrictions $\rho_0, \sigma$ on
$n$ variables, and an accumulator $\mathrm{acc}$ recorded as a list of index–value pairs
in $\bbn \times \mathrm{Bool}$; it returns a pair of lists. Then the second component of
its output has length at most $\abs{\mathrm{acc}} + 2\,\abs{\mathrm{path}}$, where
$\abs{\mathrm{acc}}$ and $\abs{\mathrm{path}}$ denote the lengths of the accumulator and
path lists.
\end{lemma}

\begin{lemma}[Auxiliary encoding length is at most twice the depth]
\lean{SwitchingLemma2.razborovEncode_aux_length_le}
\label{SwitchingLemma2.razborovEncode_aux_length_le}
\leanok
\proofsource{switching-lemma}{
  pdf-b5e074215b9e-p002-b003,
  pdf-b5e074215b9e-p002-b004
}
\proofstep
  {BoolCircuit.Lit}
  {context}
  {switching-lemma}
  {pdf-b5e074215b9e-p001-b004}
\proofstep
  {BoolCircuit.Lit.eval}
  {context}
  {switching-lemma}
  {pdf-b5e074215b9e-p001-b004}
\proofstep
  {Literal}
  {context}
  {switching-lemma}
  {pdf-b5e074215b9e-p001-b004}
\proofstep
  {DecisionTree.eval}
  {context}
  {switching-lemma}
  {pdf-b5e074215b9e-p001-b005}
\proofstep
  {Term}
  {context}
  {switching-lemma}
  {pdf-b5e074215b9e-p001-b004}
\proofstep
  {DNF}
  {context}
  {switching-lemma}
  {pdf-b5e074215b9e-p001-b004}
\proofstep
  {DecisionTree}
  {context}
  {switching-lemma}
  {pdf-b5e074215b9e-p001-b005}
\proofstep
  {DecisionTree.depth}
  {context}
  {switching-lemma}
  {pdf-b5e074215b9e-p001-b005}
\proofstep
  {DecisionTree.deepPath}
  {context}
  {switching-lemma}
  {pdf-b5e074215b9e-p002-b003}
\proofstep
  {buildFullDTree}
  {context}
  {switching-lemma}
  {pdf-b5e074215b9e-p001-b005}
\proofstep
  {buildFullDTree_depth}
  {context}
  {switching-lemma}
  {pdf-b5e074215b9e-p001-b005}
\proofstep
  {buildFullDTree_eval}
  {context}
  {switching-lemma}
  {pdf-b5e074215b9e-p001-b005}
\proofstep
  {dtDepth}
  {context}
  {switching-lemma}
  {pdf-b5e074215b9e-p001-b005}
\proofstep
  {SwitchingLemma2.Restriction}
  {context}
  {switching-lemma}
  {pdf-b5e074215b9e-p001-b007}
\proofstep
  {SwitchingLemma2.Restriction.freeVars}
  {context}
  {switching-lemma}
  {pdf-b5e074215b9e-p001-b007}
\proofstep
  {SwitchingLemma2.Restriction.numFree}
  {context}
  {switching-lemma}
  {pdf-b5e074215b9e-p001-b007}
\proofstep
  {SwitchingLemma2.Restriction.extend}
  {context}
  {switching-lemma}
  {pdf-b5e074215b9e-p001-b007}
\proofstep
  {SwitchingLemma2.Literal.fixedBy}
  {context}
  {switching-lemma}
  {pdf-b5e074215b9e-p002-b004}
\proofstep
  {SwitchingLemma2.restrictFn}
  {context}
  {switching-lemma}
  {pdf-b5e074215b9e-p001-b007}
\proofstep
  {SwitchingLemma2.IsBadRestriction}
  {context}
  {switching-lemma}
  {pdf-b5e074215b9e-p002-b001}
\proofstep
  {SwitchingLemma2.canonicalDTree}
  {context}
  {switching-lemma}
  {pdf-b5e074215b9e-p002-b002}
\proofstep
  {SwitchingLemma2.processClauseLits}
  {context}
  {switching-lemma}
  {pdf-b5e074215b9e-p002-b003}
\proofstep
  {SwitchingLemma2.razborovEncode}
  {context}
  {switching-lemma}
  {pdf-b5e074215b9e-p002-b003,
    pdf-b5e074215b9e-p002-b004}
\proofstep
  {SwitchingLemma2.processClauseLits_bound}
  {context}
  {switching-lemma}
  {pdf-b5e074215b9e-p002-b003}
\proofstep
  {SwitchingLemma2.processClauseLits_tight}
  {context}
  {switching-lemma}
  {pdf-b5e074215b9e-p002-b003}
\proofstep
  {SwitchingLemma2.encode_go_aux_length_bound}
  {context}
  {switching-lemma}
  {pdf-b5e074215b9e-p002-b004}
\uses{DNF, DNF.eval, DecisionTree.deepPath, SwitchingLemma2.IsBadRestriction, SwitchingLemma2.Restriction, SwitchingLemma2.canonicalDTree, SwitchingLemma2.encode_go_aux_length_bound, SwitchingLemma2.razborovEncode}
\difficulty{3}
Let $f$ be a DNF formula on $n$ Boolean variables, let $w, d \in \bbn$, and let $\rho$
be a restriction that is bad for $f$ at depth $d$; that is, the restriction of the
Boolean function computed by $f$ along $\rho$ has decision-tree depth strictly greater
than $d$. Then the auxiliary component of the Razborov encoding of $f$ and $\rho$ with
parameters $w$ and $d$ — the list of $(\bbn \times \mathrm{Bool})$ pairs returned as the
second entry of the encoding — has length at most $2d$.
\end{lemma}

\begin{lemma}[Stability of the restriction σ away from a clause's variables]
\lean{SwitchingLemma2.processClauseLits_sigma_stable}
\label{SwitchingLemma2.processClauseLits_sigma_stable}
\leanok
\uses{Literal, SwitchingLemma2.Restriction, SwitchingLemma2.processClauseLits}
\difficulty{4}
Fix $n$ and let $\mathrm{lits}$ be a list of literals over $\mathrm{Fin}\,n$, each
paired with a position index, let $\mathrm{path}$ be a list of variable–direction pairs,
and let $\rho_0$ and $\sigma$ be restrictions on $n$ variables. Suppose $v \in
\mathrm{Fin}\,n$ is not the variable of any literal occurring in $\mathrm{lits}$. Then
the restriction $\sigma'$ produced by processing $\mathrm{lits}$ against $\mathrm{path}$
starting from $\rho_0$ and $\sigma$ agrees with $\sigma$ at $v$, that is $\sigma'(v) =
\sigma(v)$.
\end{lemma}

\begin{lemma}[Untouched variables of a processed clause remain fixed in $\rho_0$]
\lean{SwitchingLemma2.processClauseLits_rho_stable}
\label{SwitchingLemma2.processClauseLits_rho_stable}
\leanok
\uses{Literal, SwitchingLemma2.Restriction, SwitchingLemma2.processClauseLits}
\difficulty{4}
Fix $n$ and let $\mathrm{lits}$ be a list of literals on $n$ variables, each paired with
a position index, and let $\mathrm{path}$ be a list of variable–direction pairs. Given
restrictions $\rho_0,\sigma$ on $n$ variables and a variable $v$ such that $v$ is not
the variable of any literal occurring in $\mathrm{lits}$, the updated restriction
$\rho_0'$ produced by processing the clause's free literals of $\mathrm{lits}$ against
$\mathrm{path}$ starting from $\rho_0$ and $\sigma$ agrees with $\rho_0$ at $v$; that
is, $\rho_0'(v) = \rho_0(v)$.
\end{lemma}

\begin{lemma}[Processing a clause preserves fixed variables of the restriction]
\lean{SwitchingLemma2.processClauseLits_rho_ne_none}
\label{SwitchingLemma2.processClauseLits_rho_ne_none}
\leanok
\uses{Literal, SwitchingLemma2.Restriction, SwitchingLemma2.encode_go_fst_acc, SwitchingLemma2.processClauseLits}
\difficulty{3}
Fix $n$, and let $\rho_0$ and $\sigma$ be restrictions on $n$ variables, let $\ell$ be a
list of literals each paired with its position index in a clause, and let a
decision-tree path be given. Suppose the variable $v$ is fixed under $\rho_0$, i.e.
$\rho_0(v) \neq \mathrm{none}$. Then, after processing the clause's free literals
against the path — which returns updated versions of $\rho_0$, $\sigma$, together with
the remaining path and the accumulated auxiliary data — the variable $v$ remains fixed
in the returned restriction $\rho_0$: its value there is again not $\mathrm{none}$. In
other words, this process never frees a variable that was already fixed in $\rho_0$.
\end{lemma}

\begin{lemma}[Encoder agrees with $\sigma$ at variables fixed by $\rho_0$]
\lean{SwitchingLemma2.encode_go_fst_nonfree}
\label{SwitchingLemma2.encode_go_fst_nonfree}
\leanok
\uses{DNF, Literal, SwitchingLemma2.Restriction, SwitchingLemma2.Restriction.freeVars, SwitchingLemma2.encode_go_fst_acc, SwitchingLemma2.processClauseLits_rho_ne_none, SwitchingLemma2.processClauseLits_sigma_stable}
\difficulty{5}
Fix $n$ and work with formulas on the variables $\mathrm{Fin}\,n$. Let $f$ be a DNF
formula, $w \in \bbn$ a width parameter, and $\mathrm{fuel} \in \bbn$ a bound on the
recursion depth; let $\mathrm{path}$ be a list of variable–value pairs (elements of
$\mathrm{Fin}\,n \times \mathrm{Bool}$), let $\rho_0$ and $\sigma$ be restrictions on
$n$ variables, and let $\mathrm{acc}$ be a list of index–value pairs holding the
encoding accumulated so far. The recursive core of the Razborov encoder sends this data
to a pair whose first component is again a restriction; write $\gamma$ for that
restriction. Then for every variable $v$ that is fixed by $\rho_0$ — that is, $\rho_0(v)
\neq \mathrm{none}$ — one has $\gamma(v) = \sigma(v)$.
\end{lemma}

\begin{lemma}[Accumulator decomposition for the encoder loop]
\lean{SwitchingLemma2.encode_go_acc}
\label{SwitchingLemma2.encode_go_acc}
\leanok
\uses{DNF, SwitchingLemma2.Restriction}
\difficulty{4}
Fix a DNF formula $f$ on $n$ variables — a list of terms, each a conjunction of literals
— together with a width parameter $w \in \bbn$, a fuel bound $\mathrm{fuel} \in \bbn$, a
path $P$ recording the variable–value pairs already committed (a list of elements of
$\mathrm{Fin}\,n \times \mathrm{Bool}$), two restrictions $\rho_0, \sigma :
\mathrm{Fin}\,n \to \mathrm{Option}\,\mathrm{Bool}$, and an accumulator $A$, a list of
index–value pairs in $\bbn \times \mathrm{Bool}$. The tail-recursive helper of
Razborov's encoding procedure, run from the initial accumulator $A$ on these data,
produces the same output as the corresponding run started from the empty accumulator,
with the entries of $A$ carried through unchanged; that is, the accumulator contributes
to the result only additively and does not affect the encoding computed by the procedure
itself.
\end{lemma}

\begin{lemma}[Accumulator-independence of the encoder's first output]
\lean{SwitchingLemma2.encode_go_fst_acc}
\label{SwitchingLemma2.encode_go_fst_acc}
\leanok
\uses{DNF, SwitchingLemma2.Restriction, SwitchingLemma2.encode_go_acc}
\difficulty{1}
Fix a DNF formula $f$ on $n$ variables, a width parameter $w \in \bbn$, a fuel bound
$\in \bbn$, a partial path recorded as a list of variable–value pairs $(i,b) \in
\mathrm{Fin}\,n \times \{0,1\}$, and two restrictions $\rho_0, \sigma$ on the $n$
variables. The auxiliary recursion underlying the Razborov encoding takes, in addition,
a starting accumulator: a list of index–bit pairs in $\bbn \times \{0,1\}$, and returns
a pair. Then the first component of its output does not depend on that accumulator:
running the recursion from an arbitrary accumulator yields the same first component as
running it from the empty accumulator.
\end{lemma}

\begin{lemma}[Entry processing preserves free variables]
\lean{SwitchingLemma2.processEntries_preserves_none}
\label{SwitchingLemma2.processEntries_preserves_none}
\leanok
\uses{SwitchingLemma2.Restriction, Term}
\difficulty{3}
Fix $n$ variables and consider the entry-processing step of the Razborov decoding
procedure, which takes a term $t$ (a conjunction of literals) over these variables, a
natural number $w$, two restrictions $\sigma, \rho_0 \colon \mathrm{Fin}\,n \to
\mathrm{Option}\,\mathrm{Bool}$, and a list $\mathit{aux}$ of pairs $(\mathbb{N} \times
\mathrm{Bool})$, and returns a restriction as the first component of its output. If
$\sigma$ leaves a variable $v$ free (that is, $\sigma(v) = \mathtt{none}$), then the
restriction produced by this step also leaves $v$ free.
\end{lemma}

\begin{lemma}[Decoder loop preserves free variables]
\lean{SwitchingLemma2.decode_go_preserves_none}
\label{SwitchingLemma2.decode_go_preserves_none}
\leanok
\uses{DNF, SwitchingLemma2.Restriction, SwitchingLemma2.processEntries_preserves_none}
\difficulty{4}
Fix a number of variables $n$, and let $D$ denote the recursive loop underlying
Razborov's decoder, which takes as input a DNF formula $f$ on $n$ variables, a width
parameter $w \in \bbn$, a fuel bound $\mathrm{fuel} \in \bbn$, two restrictions $\sigma$
and $\rho_0$ on the $n$ variables, and an auxiliary list of index–value pairs drawn from
$\bbn \times \mathrm{Bool}$, and returns a pair whose first component is again a
restriction on the $n$ variables. Then for every variable $v$ with $\sigma(v) =
\mathrm{none}$, the first component of $D(f, w, \mathrm{fuel}, \sigma, \rho_0,
\mathrm{aux})$ also assigns $\mathrm{none}$ to $v$.
\end{lemma}

\begin{lemma}[Independence of the path, restriction, and auxiliary outputs from $\sigma$]
\lean{SwitchingLemma2.processClauseLits_sigma_indep}
\label{SwitchingLemma2.processClauseLits_sigma_indep}
\leanok
\uses{Literal, SwitchingLemma2.Restriction, SwitchingLemma2.processClauseLits}
\difficulty{3}
Fix a list of literals, each paired with its position index within a clause, a
decision-tree path, a restriction $\rho_0$, and two arbitrary restrictions $\sigma_1$
and $\sigma_2$ on $n$ variables. Running the clause-processing procedure on this data
with $\sigma_1$ and with $\sigma_2$ produces the same three outputs in each case: the
same remaining path, the same updated restriction $\rho_0$, and the same accumulated
auxiliary list of (position index, path direction) pairs. In other words, of the
procedure's outputs only the returned $\sigma$ can depend on the $\sigma$ supplied as
input; the path, the $\rho_0$ component, and the auxiliary data are all unaffected by
that choice.
\end{lemma}

\begin{lemma}[Encoder output's second component is independent of σ]
\lean{SwitchingLemma2.encode_go_snd_sigma_indep}
\label{SwitchingLemma2.encode_go_snd_sigma_indep}
\leanok
\uses{DNF, SwitchingLemma2.Restriction, SwitchingLemma2.encode_go_acc, SwitchingLemma2.processClauseLits, SwitchingLemma2.processClauseLits_sigma_indep}
\difficulty{5}
Fix a DNF formula $f$ on $n$ variables, two natural-number parameters $w$ and $t$ (a
fuel bound), a list of decisions pairing a variable with a Boolean value, and a
restriction $\rho_0$. Run the Razborov encoding loop on this data, started with empty
auxiliary data, using a restriction $\sigma$. Then the second component of the loop's
output is the same whether one supplies the restriction $\sigma_1$ or the restriction
$\sigma_2$ in place of $\sigma$; that is, this component does not depend on the $\sigma$
argument.
\end{lemma}

\begin{lemma}[Auxiliary indices trace back to input literals]
\lean{SwitchingLemma2.processClauseLits_aux_entries_from_lits}
\label{SwitchingLemma2.processClauseLits_aux_entries_from_lits}
\leanok
\uses{Literal, SwitchingLemma2.Restriction, SwitchingLemma2.processClauseLits}
\difficulty{4}
Fix $n$, and process the free literals of a clause against a decision-tree path: take a
list $\mathrm{lits}$ of literal–index pairs, a decision-tree path, and two restrictions
$\rho_0$ and $\sigma$ on $n$ variables, and write $A$ for the accumulated list of
index–direction pairs returned as auxiliary data. Then for every pair $(i, b) \in A$,
the index $i$ equals the index component $j$ of some literal–index pair $(\ell, j) \in
\mathrm{lits}$.
\end{lemma}

\begin{lemma}[Recorded positions avoid a designated variable]
\lean{SwitchingLemma2.processClauseLits_aux_ne_nonfree}
\label{SwitchingLemma2.processClauseLits_aux_ne_nonfree}
\leanok
\uses{Literal, SwitchingLemma2.Restriction, SwitchingLemma2.processClauseLits, SwitchingLemma2.processClauseLits_aux_entries_from_lits, SwitchingLemma2.zipIdx_drop_spec, Term}
\difficulty{4}
Let $t$ be a term in $n$ variables, let $v \in \mathrm{Fin}\,n$ be a variable, and let
$\rho_0, \sigma$ be restrictions on $n$ variables. Let $\mathit{lits}$ be a list of
literals each tagged with a natural-number position, and suppose that every tagged
literal in $\mathit{lits}$ occurs in the position-indexed enumeration of $t$ (each
literal of $t$ paired with its index), and that no literal appearing in $\mathit{lits}$
has variable $v$. Consider the clause auxiliary data — the list of (position, direction)
pairs — produced by processing $\mathit{lits}$ against a path starting from $\rho_0$ and
$\sigma$. Then for every pair $(i, b)$ in this auxiliary data, whenever the literals of
$t$ from position $i$ onward form a nonempty list with leading literal $\ell$, the
variable of $\ell$ is different from $v$.
\end{lemma}

\begin{lemma}[Auxiliary entries mark variables free under the restriction]
\lean{SwitchingLemma2.processClauseLits_aux_vars_free}
\label{SwitchingLemma2.processClauseLits_aux_vars_free}
\leanok
\uses{Literal, SwitchingLemma2.Restriction, SwitchingLemma2.processClauseLits, SwitchingLemma2.processClauseLits_aux_entries_from_lits, SwitchingLemma2.zipIdx_drop_spec, Term}
\difficulty{3}
Let $t$ be a term on $n$ variables (a list of literals), let $\rho_0$ and $\sigma$ be
restrictions on $n$ variables, let $L$ be a list of literals each paired with a position
index, and let $P$ be a decision-tree path. Suppose that every pair $(\ell, i)$ in $L$
occurs in the indexed enumeration of $t$ — so that $\ell$ is precisely the literal at
position $i$ of $t$ — and that every literal appearing in $L$ has its variable free
under $\rho_0$ (its value under $\rho_0$ is $\mathrm{none}$). Then for each auxiliary
entry $(i, b)$ produced by processing $L$ against $P$ starting from $\rho_0$ and
$\sigma$, and for any literal $\ell'$ and list $r$ with $t$ from position $i$ onward
equal to $\ell' :: r$ (i.e. discarding the first $i$ literals of $t$ leaves a nonempty
list headed by $\ell'$), the variable of $\ell'$ is free under $\rho_0$.
\end{lemma}

\begin{lemma}[Fold of variable-freeing updates leaves untargeted variables fixed]
\lean{SwitchingLemma2.foldl_sigma_stable}
\label{SwitchingLemma2.foldl_sigma_stable}
\leanok
\uses{Literal, SwitchingLemma2.Restriction, Term}
\difficulty{4}
Fix a term $t$ over $n$ variables (a list of literals), a restriction $\sigma :
\mathrm{Fin}\,n \to \mathrm{Option}\,\mathrm{Bool}$, a variable $v$, and a finite list
of entries, each a pair $(k, b)$ with $k \in \bbn$ and $b$ a Boolean. Process the
entries left to right, starting from $\sigma$: for an entry $(k, b)$, delete the first
$k$ literals of $t$; if what remains is empty, leave the current restriction unchanged,
and otherwise, writing $l$ for the first remaining literal, reset the restriction to be
free (unassigned) at the variable of $l$. (The Boolean component $b$ of each entry is
never used.) Suppose that for every entry $(k, b)$ in the list, if deleting the first
$k$ literals of $t$ leaves a nonempty list whose first literal is $l$, then the variable
of $l$ is not $v$. Then the restriction obtained after processing all the entries takes
the same value at $v$ as $\sigma$ does.
\end{lemma}

\begin{lemma}[Stability of a restriction fold at an untouched variable]
\lean{SwitchingLemma2.foldl_rho_stable}
\label{SwitchingLemma2.foldl_rho_stable}
\leanok
\uses{Literal, SwitchingLemma2.Restriction, Term}
\difficulty{4}
Fix $n$, and let $t$ be a term over $n$ variables, that is a list of literals, together
with a restriction $\rho_0 : \mathrm{Fin}\,n \to \mathrm{Option}\,\mathrm{Bool}$, a
variable $v \in \mathrm{Fin}\,n$, and a finite list of entries $(k_1, b_1), \dots, (k_m,
b_m)$ with each $k_i \in \bbn$ and each $b_i$ a Boolean. Process these entries from left
to right, starting from $\rho_0$: for an entry $(k_i, b_i)$, delete the first $k_i$
literals of $t$; if what remains is nonempty with head literal $l$, update the current
restriction so that the variable of $l$ is now assigned the fixed value $b_i$, and
otherwise leave the restriction unchanged. Suppose that for every entry $(k_i, b_i)$,
whenever deleting the first $k_i$ literals of $t$ leaves a nonempty list with head
literal $l$, the variable of $l$ is not $v$. Then the restriction obtained after
processing all the entries assigns $v$ the same value as $\rho_0$ does.
\end{lemma}

\begin{lemma}[Folding never assigns a free variable]
\lean{SwitchingLemma2.foldl_sigma_preserves_none}
\label{SwitchingLemma2.foldl_sigma_preserves_none}
\leanok
\uses{SwitchingLemma2.Restriction, Term}
\difficulty{3}
Let $t$ be a term on $n$ variables (a list of literals), let $(k_1,b_1),\dots,(k_m,b_m)$
be a list of pairs in $\bbn\times\mathrm{Bool}$, and let $\rho$ be a restriction on $n$
variables. Process the entries left to right, starting from $\rho$: for the entry
$(k_i,b_i)$, delete the first $k_i$ literals of $t$; if nothing remains, leave the
current restriction unchanged, and otherwise reset the variable of the first surviving
literal to free. Then for every variable $v$ left free by $\rho$, the restriction
obtained at the end of this process still leaves $v$ free.
\end{lemma}

\begin{lemma}[Decoder fold frees a variable set during clause processing]
\lean{SwitchingLemma2.processClauseLits_foldl_sigma_none}
\label{SwitchingLemma2.processClauseLits_foldl_sigma_none}
\leanok
\uses{Literal, SwitchingLemma2.Restriction, SwitchingLemma2.foldl_sigma_preserves_none, SwitchingLemma2.processClauseLits, SwitchingLemma2.zipIdx_drop_spec, Term}
\difficulty{5}
Let $t$ be a term on $n$ variables — a list of literals, each carrying a variable and a
negation flag — and refer to the position of a literal by its index in this list. Let
$\mathrm{lits}$ be a list of (literal, index) pairs, every one of which is the literal
occurring at that index in $t$; let $\mathrm{path}$ be a list of (variable, Boolean)
pairs; and let $\rho_0$, $\sigma$, and $\sigma_{\mathrm{dec}}$ be restrictions on $n$
variables, with $v$ a distinguished variable. Feeding $\mathrm{lits}$, $\mathrm{path}$,
$\rho_0$, and $\sigma$ to the clause-processing procedure returns an updated restriction
$\rho_0'$ together with an auxiliary list $A$ of (index, direction) pairs. Suppose
$\rho_0(v)$ is free while $\rho_0'(v)$ is fixed. Then processing $A$ from left to right,
starting from $\sigma_{\mathrm{dec}}$ and, at each entry whose index component is $i$,
resetting to free the variable of the $i$-th literal of $t$ whenever $t$ has a literal
at position $i$ (and otherwise leaving the restriction unchanged), yields a restriction
that is free at $v$.
\end{lemma}

\begin{lemma}[Decoder fold recovers the clause-processing restriction]
\lean{SwitchingLemma2.processClauseLits_foldl_rho_eq}
\label{SwitchingLemma2.processClauseLits_foldl_rho_eq}
\leanok
\uses{Literal, SwitchingLemma2.Restriction, SwitchingLemma2.processClauseLits, SwitchingLemma2.zipIdx_drop_spec, Term}
\difficulty{4}
Let $t$ be a term on $n$ variables (a list of literals), and let $\mathrm{lits}$ be a
list of literal–index pairs, each of which is required to be a literal of $t$ together
with its position — that is, every pair in $\mathrm{lits}$ occurs in the list obtained
by pairing each literal of $t$ with its index. Fix a path, restrictions $\rho_0$,
$\sigma$, and $\rho_0^{\mathrm{dec}}$ on the $n$ variables, and a variable $v$ with
$\rho_0(v) = \rho_0^{\mathrm{dec}}(v)$. Processing $\mathrm{lits}$ against the path
starting from $\rho_0$ and $\sigma$ yields an updated restriction $\rho_0'$ together
with an auxiliary list of pairs $(i, b)$. Now form a restriction by folding over this
auxiliary list, left to right, starting from $\rho_0^{\mathrm{dec}}$: an entry $(i, b)$
leaves the current restriction unchanged if $t$ has no literal at position $i$, and
otherwise fixes the variable of the literal at position $i$ to the Boolean value $b$.
Then the restriction so obtained and $\rho_0'$ take the same value at $v$.
\end{lemma}

\begin{lemma}[Decoding fold recovers the processed restriction at a newly set variable]
\lean{SwitchingLemma2.processClauseLits_foldl_rho_eq_of_set}
\label{SwitchingLemma2.processClauseLits_foldl_rho_eq_of_set}
\leanok
\uses{Literal, SwitchingLemma2.Restriction, SwitchingLemma2.processClauseLits, SwitchingLemma2.processClauseLits_foldl_rho_eq, SwitchingLemma2.processClauseLits_no_target_of_rho_none, SwitchingLemma2.zipIdx_drop_spec, Term}
\difficulty{5}
Let $t$ be a term on $n$ variables, and let $L$ be a list of literal–index pairs each of
which occurs at its recorded position in $t$ (that is, every pair in $L$ lies in the
enumeration of $t$ that tags each literal with its index). Fix restrictions $\rho_0$,
$\sigma$, and $\rho_0^{\mathrm{dec}}$ on $n$ variables, a decision-tree path $P$, and a
variable $v$ that $\rho_0$ leaves free. Running the clause-processing routine on $L$,
$P$, $\rho_0$, $\sigma$ produces an updated restriction $\rho_0'$ together with
auxiliary data $A$, a list of (position index, direction) pairs, and assume that
$\rho_0'$ assigns a value at $v$. Consider the restriction obtained by folding $A$ over
the starting restriction $\rho_0^{\mathrm{dec}}$, where each pair $(i, b)$ leaves the
current restriction unchanged if $t$ has no literal at position $i$, and otherwise sets
the variable of the $i$-th literal of $t$ to the value $b$. Then this folded restriction
agrees with $\rho_0'$ at $v$; that is, its value at $v$ equals $\rho_0'(v)$.
\end{lemma}

\begin{lemma}[Still-free variables avoid the clause's literals]
\lean{SwitchingLemma2.processClauseLits_no_target_of_rho_none}
\label{SwitchingLemma2.processClauseLits_no_target_of_rho_none}
\leanok
\uses{Literal, SwitchingLemma2.Restriction, SwitchingLemma2.processClauseLits, SwitchingLemma2.processClauseLits_rho_ne_none}
\difficulty{5}
Fix a number $n$ of variables, and let $\rho_0$ and $\sigma$ be restrictions on them.
Suppose the free literals of a clause are given as a list, each literal paired with its
position index, and let a decision-tree path be a list of variable–direction pairs whose
length is at least the number of literals. Running the clause-processing operation on
this list, this path, $\rho_0$, and $\sigma$ returns an updated restriction $\rho_0'$ in
place of $\rho_0$. If a variable $v$ is free under $\rho_0$ and remains free under
$\rho_0'$, then no literal on the list has $v$ as its underlying variable.
\end{lemma}

\begin{lemma}[Updated restriction at a variable depends only on its starting value]
\lean{SwitchingLemma2.processClauseLits_sigma_at_v}
\label{SwitchingLemma2.processClauseLits_sigma_at_v}
\leanok
\uses{Literal, SwitchingLemma2.Restriction, SwitchingLemma2.processClauseLits}
\difficulty{3}
Consider the routine that processes one clause's free literals — a list of literals,
each tagged with its position index in the clause — against a canonical decision-tree
path, starting from a restriction $\rho_0$ and a restriction $\sigma$, and returning
(among other data) an updated restriction $\sigma'$. Fix the list of literals, the path,
the restriction $\rho_0$, and a variable $v$. If two starting restrictions $\sigma_1$
and $\sigma_2$ agree at $v$, that is $\sigma_1(v) = \sigma_2(v)$, then the updated
restrictions produced from $\sigma_1$ and from $\sigma_2$ also agree at $v$;
equivalently, the value at $v$ of the updated $\sigma$-restriction depends only on the
starting value $\sigma(v)$.
\end{lemma}

\begin{lemma}[Processed assignment unchanged where the restriction stays free]
\lean{SwitchingLemma2.processClauseLits_sigma_none_of_rho_none}
\label{SwitchingLemma2.processClauseLits_sigma_none_of_rho_none}
\leanok
\uses{Literal, SwitchingLemma2.Restriction, SwitchingLemma2.processClauseLits, SwitchingLemma2.processClauseLits_rho_ne_none}
\difficulty{4}
Fix $n$, and let $\rho_0$ and $\sigma$ be restrictions on $n$ variables. Run the
clause-processing procedure on a list of literals (each paired with its position index)
together with a decision-tree path, starting from $\rho_0$ and $\sigma$, and write
$\rho_0'$ and $\sigma'$ for the two restrictions it returns. If a variable $v$ is free
under the input, $\rho_0(v) = \mathrm{none}$, and remains free under the output,
$\rho_0'(v) = \mathrm{none}$, then the procedure has not altered $\sigma$ at $v$ either:
$\sigma'(v) = \sigma(v)$.
\end{lemma}

\begin{lemma}[Free-variable output of the Razborov encoder is independent of its starting restriction]
\lean{SwitchingLemma2.encode_go_fst_sigma_indep_at_free}
\label{SwitchingLemma2.encode_go_fst_sigma_indep_at_free}
\leanok
\uses{DNF, SwitchingLemma2.Restriction, SwitchingLemma2.encode_go_fst_acc, SwitchingLemma2.encode_go_fst_nonfree, SwitchingLemma2.processClauseLits, SwitchingLemma2.processClauseLits_sigma_at_v, SwitchingLemma2.processClauseLits_sigma_indep, SwitchingLemma2.processClauseLits_sigma_none_of_rho_none}
\difficulty{5}
Fix a DNF formula $f$ on $n$ variables, a width $w$, a fuel bound, a decision path (a
list of variable–direction pairs), and a base restriction $\rho_0$. The recursive worker
of the Razborov encoding, run on these inputs together with a running restriction
$\sigma$ and empty auxiliary data, returns among its outputs a restriction; write
$R(\sigma)$ for this first output component. Let $v$ be a variable that is left free by
$\rho_0$ and by two running restrictions $\sigma_1$ and $\sigma_2$, meaning $\rho_0(v)$,
$\sigma_1(v)$, and $\sigma_2(v)$ are all undefined. Then the two runs agree at $v$:
$R(\sigma_1)(v) = R(\sigma_2)(v)$.
\end{lemma}

\begin{lemma}[Decoder replay of a clause's processed free literals]
\lean{SwitchingLemma2.processEntries_of_processClauseLits}
\proofsource{switching-lemma}{pdf-b5e074215b9e-p002-b004}
\label{SwitchingLemma2.processEntries_of_processClauseLits}
\leanok
\proofstep
  {BoolCircuit.Lit}
  {context}
  {switching-lemma}
  {pdf-b5e074215b9e-p001-b004}
\proofstep
  {BoolCircuit.Lit.eval}
  {context}
  {switching-lemma}
  {pdf-b5e074215b9e-p001-b004}
\proofstep
  {Literal}
  {context}
  {switching-lemma}
  {pdf-b5e074215b9e-p001-b004}
\proofstep
  {Term}
  {context}
  {switching-lemma}
  {pdf-b5e074215b9e-p001-b004}
\proofstep
  {DNF}
  {context}
  {switching-lemma}
  {pdf-b5e074215b9e-p001-b004}
\proofstep
  {SwitchingLemma2.Restriction}
  {context}
  {switching-lemma}
  {pdf-b5e074215b9e-p001-b007}
\proofstep
  {SwitchingLemma2.Literal.fixedBy}
  {context}
  {switching-lemma}
  {pdf-b5e074215b9e-p002-b004}
\proofstep
  {SwitchingLemma2.zipIdx_drop_spec}
  {context}
  {switching-lemma}
  {pdf-b5e074215b9e-p002-b003}
\proofstep
  {SwitchingLemma2.processClauseLits}
  {context}
  {switching-lemma}
  {pdf-b5e074215b9e-p002-b003}
\proofstep
  {SwitchingLemma2.razborovDecode}
  {context}
  {switching-lemma}
  {pdf-b5e074215b9e-p002-b004}
\uses{Literal, SwitchingLemma2.Restriction, SwitchingLemma2.processClauseLits, SwitchingLemma2.zipIdx_drop_spec, Term}
\difficulty{5}
Work over $n$ Boolean variables, where a *literal* is a variable together with a
negation flag, a *term* is a list of literals (a conjunction), and a *restriction* $\rho
: \mathrm{Fin}\,n \to \mathrm{Option}\,\mathrm{Bool}$ fixes some variables and leaves
the rest free. Let $t$ be a term whose length is at most a bound $w$, let
$\mathit{lits}$ be a list of literals each paired with a position index, every such pair
occurring among the indexed entries of $t$, and fix a canonical decision-tree path
together with restrictions $\rho_0^{\mathrm{enc}}, \sigma^{\mathrm{enc}},
\sigma^{\mathrm{dec}}, \rho_0^{\mathrm{dec}}$ and a remainder list. Then, running the
decoder's entry-processing routine on the auxiliary data emitted when this clause's free
literals are consumed against the path — followed by the termination marker $(w,
\mathrm{false})$ and the remainder — yields the same $\sigma$- and $\rho_0$-restrictions
as folding the corresponding per-entry updates directly over that auxiliary data, with
the remainder left unconsumed.
\end{lemma}

\begin{lemma}[Processing free literals preserves a variable's non-negating assignment]
\lean{SwitchingLemma2.processClauseLits_sigma_ne_neg}
\label{SwitchingLemma2.processClauseLits_sigma_ne_neg}
\leanok
\uses{Literal, SwitchingLemma2.Restriction, SwitchingLemma2.processClauseLits}
\difficulty{4}
Let $\ell$ be a literal with variable $v$ and negation flag $b$, let $\rho_0$ and
$\sigma$ be restrictions on $n$ variables, and let $\mathit{lits}$ be a list of literals
(each tagged with its position index in the clause) run against a decision-tree path.
Suppose that every literal in $\mathit{lits}$ whose variable is $v$ is equal to $\ell$,
and that $\sigma$ does not already assign $v$ the fixed value $b$, i.e. $\sigma(v) \neq
\mathrm{some}\,b$. Then the restriction $\sigma'$ obtained as the $\sigma$-component of
processing $\mathit{lits}$ against the path starting from $\rho_0$ and $\sigma$ still
does not assign $v$ the value $b$, i.e. $\sigma'(v) \neq \mathrm{some}\,b$.
\end{lemma}

\begin{lemma}[Empty remaining path when a clause literal's variable stays free]
\lean{SwitchingLemma2.processClauseLits_path_nil_of_rho_none_and_mem}
\label{SwitchingLemma2.processClauseLits_path_nil_of_rho_none_and_mem}
\leanok
\uses{Literal, SwitchingLemma2.Restriction, SwitchingLemma2.processClauseLits, SwitchingLemma2.processClauseLits_rho_ne_none}
\difficulty{5}
Fix $n$, a list $\mathit{lits}$ whose entries are literals paired with clause positions,
a path $\mathit{path}$ presented as a list of variable–direction pairs, and two
restrictions $\rho_0,\sigma$ on $n$ variables. Suppose the pair $(l,\mathit{idx})$
occurs in $\mathit{lits}$, that every entry of $\mathit{lits}$ whose variable equals
$\mathrm{var}(l)$ is equal to $l$ itself, and that $\rho_0$ leaves $\mathrm{var}(l)$
free. Processing $\mathit{lits}$ against $\mathit{path}$ from the starting restrictions
$\rho_0$ and $\sigma$ yields an updated restriction $\rho_0'$, an updated restriction
$\sigma'$, a remaining path, and accumulated auxiliary data. If $\rho_0'$ still leaves
$\mathrm{var}(l)$ free, then the remaining path yielded is empty.
\end{lemma}

\begin{lemma}[A clause literal's variable is fixed in the output restriction]
\lean{SwitchingLemma2.processClauseLits_rho_ne_none_of_mem}
\label{SwitchingLemma2.processClauseLits_rho_ne_none_of_mem}
\leanok
\uses{Literal, SwitchingLemma2.Restriction, SwitchingLemma2.processClauseLits, SwitchingLemma2.processClauseLits_rho_ne_none}
\difficulty{4}
Let $\rho_0$ and $\sigma$ be restrictions on $n$ variables, let a list of literals be
given (each paired with its position index inside its clause), and let a decision-tree
path be given as a list of variable–direction pairs, with at least as many path entries
as literals in the list. Process the literals against the path starting from $\rho_0$
and $\sigma$, and let $\rho_0'$ be the updated first restriction that this returns. If
some literal in the list has underlying variable $v$, then $v$ is not free in $\rho_0'$;
that is, $\rho_0'(v)$ is a fixed Boolean value rather than the free marker.
\end{lemma}

\begin{lemma}[Round-trip base case for the decoder]
\lean{SwitchingLemma2.roundtrip_base}
\label{SwitchingLemma2.roundtrip_base}
\leanok
\uses{DNF, SwitchingLemma2.Restriction}
\difficulty{3}
Fix $n$, and let $f$ be a DNF formula on $n$ variables and $w$ a natural number. Let
$\rho_0$, $\sigma$, $\sigma_{\mathrm{dec}}$, and $\rho_0^{\mathrm{dec}}$ be restrictions
on $n$ variables (each variable either assigned a Boolean value or left free), and let
$d$ be a fuel bound. Suppose that $\sigma$ leaves free every variable that $\rho_0$
leaves free, and that $\sigma_{\mathrm{dec}}$ agrees with $\sigma$ on every variable —
both those that $\rho_0$ leaves free and those that $\rho_0$ fixes. Then running the
Razborov decoding loop on $f$, with width $w$, fuel $d$, working restrictions
$\sigma_{\mathrm{dec}}$ and $\rho_0^{\mathrm{dec}}$, and the empty list of emitted
advice, yields an output whose first component is exactly $\sigma$.
\end{lemma}

%%% added by the blueprint pipeline
\section{Additional declarations}

\begin{lemma}[Encoder does not kill the first surviving clause]
\lean{SwitchingLemma2.encode_go_not_kills_first_clause}
\proofsource{switching-lemma}{
  pdf-b5e074215b9e-p002-b003,
  pdf-b5e074215b9e-p002-b004
}
\label{SwitchingLemma2.encode_go_not_kills_first_clause}
\leanok
\proofstep
  {BoolCircuit.Lit}
  {context}
  {switching-lemma}
  {pdf-b5e074215b9e-p001-b004}
\proofstep
  {BoolCircuit.Lit.eval}
  {context}
  {switching-lemma}
  {pdf-b5e074215b9e-p001-b004}
\proofstep
  {Literal}
  {context}
  {switching-lemma}
  {pdf-b5e074215b9e-p001-b004}
\proofstep
  {Term}
  {context}
  {switching-lemma}
  {pdf-b5e074215b9e-p001-b004}
\proofstep
  {DNF}
  {context}
  {switching-lemma}
  {pdf-b5e074215b9e-p001-b004}
\proofstep
  {DecisionTree}
  {context}
  {switching-lemma}
  {pdf-b5e074215b9e-p001-b005}
\proofstep
  {DecisionTree.depth}
  {context}
  {switching-lemma}
  {pdf-b5e074215b9e-p001-b005}
\proofstep
  {DecisionTree.deepPath}
  {context}
  {switching-lemma}
  {pdf-b5e074215b9e-p002-b003,
    pdf-b5e074215b9e-p001-b005}
\proofstep
  {SwitchingLemma2.Restriction}
  {context}
  {switching-lemma}
  {pdf-b5e074215b9e-p001-b007}
\proofstep
  {SwitchingLemma2.Restriction.freeVars}
  {context}
  {switching-lemma}
  {pdf-b5e074215b9e-p001-b007}
\proofstep
  {SwitchingLemma2.Restriction.numFree}
  {context}
  {switching-lemma}
  {pdf-b5e074215b9e-p001-b007}
\proofstep
  {SwitchingLemma2.Literal.fixedBy}
  {context}
  {switching-lemma}
  {pdf-b5e074215b9e-p002-b004}
\proofstep
  {SwitchingLemma2.Term.killedBy}
  {context}
  {switching-lemma}
  {pdf-b5e074215b9e-p002-b002}
\proofstep
  {SwitchingLemma2.canonicalDTree}
  {context}
  {switching-lemma}
  {pdf-b5e074215b9e-p002-b002}
\proofstep
  {SwitchingLemma2.processClauseLits}
  {context}
  {switching-lemma}
  {pdf-b5e074215b9e-p002-b003}
\proofstep
  {SwitchingLemma2.razborovEncode}
  {context}
  {switching-lemma}
  {pdf-b5e074215b9e-p002-b003}
\proofstep
  {SwitchingLemma2.processClauseLits_sigma_stable}
  {context}
  {switching-lemma}
  {pdf-b5e074215b9e-p002-b004}
\proofstep
  {SwitchingLemma2.processClauseLits_rho_ne_none}
  {context}
  {switching-lemma}
  {pdf-b5e074215b9e-p002-b003}
\proofstep
  {SwitchingLemma2.encode_go_fst_nonfree}
  {context}
  {switching-lemma}
  {pdf-b5e074215b9e-p002-b004}
\proofstep
  {SwitchingLemma2.encode_go_acc}
  {context}
  {switching-lemma}
  {pdf-b5e074215b9e-p002-b003}
\proofstep
  {SwitchingLemma2.encode_go_fst_acc}
  {context}
  {switching-lemma}
  {pdf-b5e074215b9e-p002-b003}
\proofstep
  {SwitchingLemma2.processClauseLits_sigma_ne_neg}
  {context}
  {switching-lemma}
  {pdf-b5e074215b9e-p002-b004}
\proofstep
  {SwitchingLemma2.processClauseLits_path_nil_of_rho_none_and_mem}
  {context}
  {switching-lemma}
  {pdf-b5e074215b9e-p002-b004}
\uses{DNF, Literal, SwitchingLemma2.Restriction, SwitchingLemma2.Term.killedBy, SwitchingLemma2.encode_go_fst_acc, SwitchingLemma2.encode_go_fst_nonfree, SwitchingLemma2.processClauseLits, SwitchingLemma2.processClauseLits_path_nil_of_rho_none_and_mem, SwitchingLemma2.processClauseLits_sigma_ne_neg, Term}
\difficulty{6}
Let $f$ be a DNF formula on $n$ variables in which no term contains two distinct
literals over the same variable, and let $w$ be a natural-number width parameter. Fix
two restrictions $\rho_0$ and $\sigma$ on the $n$ variables such that $\sigma$ leaves a
variable free whenever $\rho_0$ does (that is, $\rho_0(v)$ undefined implies $\sigma(v)$
undefined), and let $t$ be the first term of $f$, in list order, that is not killed by
$\rho_0$. Run the recursion at the heart of the Razborov encoding on $f$ and $w$ — with
any fuel budget, any decision path, initial data $\rho_0$ and $\sigma$, and empty
accumulated data — and write $\rho$ for the restriction it returns as its first
component. Then for every literal $l$ of $t$ whose variable is left free by $\rho_0$,
the output restriction $\rho$ does not kill $l$: it does not fix the variable of $l$ to
the value that would make $l$ false.
\end{lemma}
